\subsection{Delta secondary}
\label{sec:na-delta}

A 480\,V delta primary with a 240\,V delta secondary has no phase shift.
The secondary is operated in one of three ways, Figure~\ref{fig:na-delta}:
ungrounded, with one corner grounded, or with a centre tap on one winding.

\begin{figure}[!ht]
  \centering
  % generated by north-american-practice/fig/connection_figures.py -- do not edit
\begin{tikzpicture}[x=0.27mm, y=-0.27mm, line cap=round, line join=round]
  \draw[hv] (38.0,130.0) -- (90.0,40.0);
  \node[font=\footnotesize, text=atInfo, anchor=center, inner sep=0.5pt] at (51.0,77.5) {480};
  \draw[hv] (90.0,40.0) -- (142.0,130.0);
  \node[font=\footnotesize, text=atInfo, anchor=center, inner sep=0.5pt] at (129.0,77.5) {480};
  \draw[hv] (142.0,130.0) -- (38.0,130.0);
  \node[font=\footnotesize, text=atInfo, anchor=center, inner sep=0.5pt] at (90.0,145.0) {480};
  \fill[atInfo] (38.0,130.0) circle (3.0);
  \node[font=\footnotesize, text=atInfo, anchor=center, inner sep=0.5pt] at (26.8,136.5) {H1};
  \fill[atInfo] (90.0,40.0) circle (3.0);
  \node[font=\footnotesize, text=atInfo, anchor=center, inner sep=0.5pt] at (90.0,27.0) {H2};
  \fill[atInfo] (142.0,130.0) circle (3.0);
  \node[font=\footnotesize, text=atInfo, anchor=center, inner sep=0.5pt] at (153.2,136.5) {H3};
  \draw[lv] (266.0,130.0) -- (318.0,40.0);
  \node[font=\footnotesize, text=atAccentFg, anchor=center, inner sep=0.5pt] at (279.0,77.5) {240};
  \draw[lv] (318.0,40.0) -- (370.0,130.0);
  \node[font=\footnotesize, text=atAccentFg, anchor=center, inner sep=0.5pt] at (357.0,77.5) {240};
  \draw[lv] (370.0,130.0) -- (266.0,130.0);
  \node[font=\footnotesize, text=atAccentFg, anchor=center, inner sep=0.5pt] at (318.0,145.0) {240};
  \fill[atAccentFg] (266.0,130.0) circle (3.0);
  \node[font=\footnotesize, text=atAccentFg, anchor=center, inner sep=0.5pt] at (254.8,136.5) {X1};
  \fill[atAccentFg] (318.0,40.0) circle (3.0);
  \node[font=\footnotesize, text=atAccentFg, anchor=center, inner sep=0.5pt] at (318.0,27.0) {X2};
  \fill[atAccentFg] (370.0,130.0) circle (3.0);
  \node[font=\footnotesize, text=atAccentFg, anchor=center, inner sep=0.5pt] at (381.2,136.5) {X3};
  \draw[arr] (172.0,104.0) -- (204.0,104.0);
  \node[font=\footnotesize, text=atFgMuted, anchor=center, inner sep=0.5pt] at (188.0,122.0) {0\textdegree};
  \node[font=\footnotesize, text=atFgMuted, anchor=center, inner sep=0.5pt] at (90.0,196.0) {480\,V $\Delta$};
  \node[font=\footnotesize, text=atFgMuted, anchor=center, inner sep=0.5pt] at (318.0,196.0) {240\,V $\Delta$, ungrounded};
  \node[font=\small\bfseries, text=atFg, anchor=west, inner sep=0.5pt] at (-6.0,18.0) {(a)};
  \draw[hv] (38.0,345.0) -- (90.0,255.0);
  \node[font=\footnotesize, text=atInfo, anchor=center, inner sep=0.5pt] at (51.0,292.5) {480};
  \draw[hv] (90.0,255.0) -- (142.0,345.0);
  \node[font=\footnotesize, text=atInfo, anchor=center, inner sep=0.5pt] at (129.0,292.5) {480};
  \draw[hv] (142.0,345.0) -- (38.0,345.0);
  \node[font=\footnotesize, text=atInfo, anchor=center, inner sep=0.5pt] at (90.0,360.0) {480};
  \fill[atInfo] (38.0,345.0) circle (3.0);
  \node[font=\footnotesize, text=atInfo, anchor=center, inner sep=0.5pt] at (26.8,351.5) {H1};
  \fill[atInfo] (90.0,255.0) circle (3.0);
  \node[font=\footnotesize, text=atInfo, anchor=center, inner sep=0.5pt] at (90.0,242.0) {H2};
  \fill[atInfo] (142.0,345.0) circle (3.0);
  \node[font=\footnotesize, text=atInfo, anchor=center, inner sep=0.5pt] at (153.2,351.5) {H3};
  \draw[lv] (266.0,345.0) -- (318.0,255.0);
  \node[font=\footnotesize, text=atAccentFg, anchor=center, inner sep=0.5pt] at (279.0,292.5) {240};
  \draw[lv] (318.0,255.0) -- (370.0,345.0);
  \node[font=\footnotesize, text=atAccentFg, anchor=center, inner sep=0.5pt] at (357.0,292.5) {240};
  \draw[lv] (370.0,345.0) -- (266.0,345.0);
  \node[font=\footnotesize, text=atAccentFg, anchor=center, inner sep=0.5pt] at (318.0,360.0) {240};
  \draw[gnd] (370.0,348.0) -- (370.0,358.0);
  \draw[gnd] (363.0,358.0) -- (377.0,358.0);
  \draw[gnd] (365.5,362.0) -- (374.5,362.0);
  \draw[gnd] (368.0,366.0) -- (372.0,366.0);
  \fill[atAccentFg] (266.0,345.0) circle (3.0);
  \node[font=\footnotesize, text=atAccentFg, anchor=center, inner sep=0.5pt] at (254.8,351.5) {X1};
  \fill[atAccentFg] (318.0,255.0) circle (3.0);
  \node[font=\footnotesize, text=atAccentFg, anchor=center, inner sep=0.5pt] at (318.0,242.0) {X2};
  \fill[atAccentFg] (370.0,345.0) circle (3.0);
  \node[font=\footnotesize, text=atAccentFg, anchor=center, inner sep=0.5pt] at (381.2,351.5) {X3};
  \draw[arr] (172.0,319.0) -- (204.0,319.0);
  \node[font=\footnotesize, text=atFgMuted, anchor=center, inner sep=0.5pt] at (188.0,337.0) {0\textdegree};
  \node[font=\footnotesize, text=atFgMuted, anchor=center, inner sep=0.5pt] at (90.0,411.0) {480\,V $\Delta$};
  \node[font=\footnotesize, text=atFgMuted, anchor=center, inner sep=0.5pt] at (318.0,411.0) {240\,V $\Delta$, X3 grounded};
  \node[font=\small\bfseries, text=atFg, anchor=west, inner sep=0.5pt] at (-6.0,233.0) {(b)};
  \draw[hv] (38.0,560.0) -- (90.0,470.0);
  \node[font=\footnotesize, text=atInfo, anchor=center, inner sep=0.5pt] at (51.0,507.5) {480};
  \draw[hv] (90.0,470.0) -- (142.0,560.0);
  \node[font=\footnotesize, text=atInfo, anchor=center, inner sep=0.5pt] at (129.0,507.5) {480};
  \draw[hv] (142.0,560.0) -- (38.0,560.0);
  \node[font=\footnotesize, text=atInfo, anchor=center, inner sep=0.5pt] at (90.0,575.0) {480};
  \fill[atInfo] (38.0,560.0) circle (3.0);
  \node[font=\footnotesize, text=atInfo, anchor=center, inner sep=0.5pt] at (26.8,566.5) {H1};
  \fill[atInfo] (90.0,470.0) circle (3.0);
  \node[font=\footnotesize, text=atInfo, anchor=center, inner sep=0.5pt] at (90.0,457.0) {H2};
  \fill[atInfo] (142.0,560.0) circle (3.0);
  \node[font=\footnotesize, text=atInfo, anchor=center, inner sep=0.5pt] at (153.2,566.5) {H3};
  \draw[lv] (266.0,560.0) -- (318.0,470.0);
  \node[font=\footnotesize, text=atAccentFg, anchor=center, inner sep=0.5pt] at (279.0,507.5) {240};
  \draw[lv] (318.0,470.0) -- (370.0,560.0);
  \node[font=\footnotesize, text=atAccentFg, anchor=center, inner sep=0.5pt] at (357.0,507.5) {240};
  \draw[lv] (370.0,560.0) -- (266.0,560.0);
  \draw[high] (318.0,560.0) -- (318.0,470.0);
  \node[font=\footnotesize, text=atDanger, anchor=east, inner sep=0.5pt] at (313.0,538.0) {208};
  \node[font=\footnotesize, text=atAccentFg, anchor=center, inner sep=0.5pt] at (291.0,552.0) {120};
  \node[font=\footnotesize, text=atAccentFg, anchor=center, inner sep=0.5pt] at (345.0,552.0) {120};
  \fill[atFg] (318.0,560.0) circle (3.0);
  \draw[gnd] (318.0,563.0) -- (318.0,573.0);
  \draw[gnd] (311.0,573.0) -- (325.0,573.0);
  \draw[gnd] (313.5,577.0) -- (322.5,577.0);
  \draw[gnd] (316.0,581.0) -- (320.0,581.0);
  \node[font=\footnotesize, text=atFg, anchor=west, inner sep=0.5pt] at (327.0,576.0) {X0};
  \fill[atAccentFg] (266.0,560.0) circle (3.0);
  \node[font=\footnotesize, text=atAccentFg, anchor=center, inner sep=0.5pt] at (254.8,566.5) {X1};
  \fill[atAccentFg] (318.0,470.0) circle (3.0);
  \node[font=\footnotesize, text=atAccentFg, anchor=center, inner sep=0.5pt] at (318.0,457.0) {X2};
  \fill[atAccentFg] (370.0,560.0) circle (3.0);
  \node[font=\footnotesize, text=atAccentFg, anchor=center, inner sep=0.5pt] at (381.2,566.5) {X3};
  \draw[arr] (172.0,534.0) -- (204.0,534.0);
  \node[font=\footnotesize, text=atFgMuted, anchor=center, inner sep=0.5pt] at (188.0,552.0) {0\textdegree};
  \node[font=\footnotesize, text=atFgMuted, anchor=center, inner sep=0.5pt] at (90.0,626.0) {480\,V $\Delta$};
  \node[font=\footnotesize, text=atFgMuted, anchor=center, inner sep=0.5pt] at (318.0,626.0) {240/120\,V $\Delta$, centre tap};
  \node[font=\small\bfseries, text=atFg, anchor=west, inner sep=0.5pt] at (-6.0,448.0) {(c)};
\end{tikzpicture}

  \caption{480\,V delta -- 240\,V delta: (a) ungrounded, (b) corner-grounded
  at X3, (c) centre tap on the winding X3--X1 with the high leg X2.}
  \label{fig:na-delta}
\end{figure}

\subsubsection{Ungrounded delta}
No conductor is grounded, and a first ground fault does not interrupt the
supply. Ungrounded systems of 120 to 1000\,V require ground detectors at the
source~[3, 250.21(B)], and the fault has to be located and cleared before
a second fault on another phase produces a phase-to-phase short circuit.
An intermittent arcing ground fault can raise the phase-to-ground voltage
above the line-to-line value~[4].

\subsubsection{Corner-grounded delta}
One phase, X3 in Figure~\ref{fig:na-delta}(b), is solidly grounded. The
other two phases are at 240\,V to ground. The grounded phase is identified
white or gray~[3, 200.6] and carries no overcurrent device unless that
device opens all conductors of the circuit together~[3, 240.22].

\subsubsection{Centre-tapped delta (high-leg delta)}
One winding, X3--X1, has a centre tap that is brought out as the neutral X0
and grounded, Figure~\ref{fig:na-delta}(c). The secondary supplies
240\,V three-phase and 120/240\,V single-phase from X1 and X3. The third
phase X2 is at
\begin{equation}
  U_{x2\_x0} = \frac{\sqrt{3}}{2} \cdot U_{ll} = \frac{\sqrt{3}}{2} \cdot 240\,\mathrm{V} = 207.8\,\mathrm{V}
  \label{eq:na-highleg}
\end{equation}
to neutral. This high leg is marked orange~[3, 110.15] and occupies the B
position in panelboards~[3, 408.3(E)]; it does not supply 120\,V loads.
The centre-tapped winding carries the single-phase load in addition to its
share of the three-phase load, so the permissible single-phase load is set
by the rating of that winding and is stated in the manufacturer's data.
\FloatBarrier
